% ══════════════════════════════════════════════════════════════════════════════
%  EXPERIMENT 2 — KEY THE SCHOOLYARD                     (Lesson 4, T3 Week 1)
% ══════════════════════════════════════════════════════════════════════════════
\experimentdiv{Key the Schoolyard}{Observable features \textbullet{} Building a dichotomous key \textbullet{} Testing it on someone else}{4}\label{sec:2}
\lessoncodes{VC2S8U01 \textbullet{} VC2S8I02, I04, I06, I08}

\vocabhead
\begin{ctab}
\begin{tabular}{|SW{4.2cm}|W{11.6cm}|}
  \hline\hd & \thead{Word} & \thead{Definition (copy from the slide)} \\ \hline
  & \scaf{\textbf{Dichotomous key}} & \scaf{A tool that identifies an organism through a series of questions, each with two possible answers.} \\ \hline
  & \textbf{Feature} & \\ \hline
  & \textbf{Observable} & \\ \hline
  & \textbf{Specimen} & \\ \hline
  & \textbf{Margin} & \\ \hline
\end{tabular}
\end{ctab}

\begin{lipanel}
\inquiryli{Build a model (a key) from careful observations, and test whether someone else can use it.}
\understandli{Classify organisms by their observable features.}
\sctitle
\begin{itemize}[leftmargin=1.4em, itemsep=0pt, topsep=1pt]
  \item I can write yes/no questions about features anyone can see.
  \item I can draw my key as a branching diagram.
  \item I can find the questions in my key that confused someone, and fix them.
\end{itemize}
\end{lipanel}

\begin{skillsfocus}
\focus{PC5}{Collect specimens without harming living things.}
\focus{PC6}{Observe each specimen closely.}
\focus{PA4}{Build a key: a model for sorting.}
\focus{CO2}{Draw the key as a branching diagram.}
\focus{EV3}{Find and fix the questions that caused errors.}
\end{skillsfocus}

% ─── TECHNIQUE ────────────────────────────────────────────────────────────────
\exppart{Technique spotlight}{Blue}

\begin{technique}{Writing a good question for a key}{PA4}
Each question splits the specimens into two groups, using a feature anyone can \textbf{see or measure}:
\begin{itemize}[leftmargin=1.4em, itemsep=1pt, topsep=3pt]
  \item \textbf{Good:} \emph{Is the leaf longer than 5\,cm?} \quad \emph{Does the leaf margin have teeth?} \quad \emph{Is the leaf split into smaller leaflets?}
  \item \textbf{Poor:} \emph{Is the leaf pretty?} (opinion) \quad \emph{Is it big?} (big compared with what?) \quad \emph{Is it a gum leaf?} (that is the answer, not a feature)
\end{itemize}
Split the specimens roughly in half each time and the key stays short.
\end{technique}

\begin{workedex}{A key for four leaves}
\begin{center}
\begin{tikzpicture}[font=\small, node distance=6mm,
  q/.style={draw=Ink, line width=0.8pt, rounded corners=1.5pt, align=center, inner sep=3pt, text width=3.6cm},
  a/.style={draw=Ink, line width=1.4pt, align=center, inner sep=3pt, font=\headingfont\bfseries\small}]
  \node[q] (q1) {Is the leaf split into leaflets?};
  \node[q, below left=8mm and -6mm of q1] (q2) {Does the margin have teeth?};
  \node[a, below right=8mm and 4mm of q1] (w) {Wattle};
  \node[a, below left=8mm and -4mm of q2] (r) {Rose};
  \node[a, below right=8mm and -4mm of q2] (g) {Gum};
  \draw[-{Stealth[length=4mm]}, Ink] (q1) -- node[left, font=\scriptsize] {no} (q2);
  \draw[-{Stealth[length=4mm]}, Ink] (q1) -- node[right, font=\scriptsize] {yes} (w);
  \draw[-{Stealth[length=4mm]}, Ink] (q2) -- node[left, font=\scriptsize] {yes} (r);
  \draw[-{Stealth[length=4mm]}, Ink] (q2) -- node[right, font=\scriptsize] {no} (g);
\end{tikzpicture}
\end{center}
\end{workedex}

\begin{wedo}{Fix the question}
A student's key asks \emph{Is the leaf a nice shape?} Explain what is wrong with it and write a better question.
\answerbox{3.5cm}
\end{wedo}

% ─── PLAN AND SAFETY / METHOD ────────────────────────────────────────────────
\exppart{Plan, safety and method}{Teal}

\textbf{Aim:} To build a dichotomous key that identifies eight schoolyard leaves, and test whether another group can use it.

\begin{infobox}{Collecting responsibly}
Collect \textbf{fallen} leaves wherever you can, and take only one leaf from a living plant. Leave animals where you find them: if you want a minibeast in your key, draw or photograph it and put it back. Scientists collect only what they need.
\end{infobox}

\needspace{10\baselineskip}
\task Complete the risk assessment. \skr{PC3}
\begin{ctab}
\begin{tabular}{|SW{4.0cm}|W{5.2cm}|W{5.4cm}|}
  \hline\hd & \thead{Risk (what could go wrong)} & \thead{Prevention} & \thead{Responding (if it happens)} \\ \hline
  & \scaf{A plant irritates the skin} & \scaf{Do not touch plants with milky sap or hairs; wash hands afterwards} & \scaf{Wash the skin with water; tell the teacher} \\ \hline
  & A bee or ant sting & & \\ \hline
  & Sunburn & & \\ \hline
\end{tabular}
\end{ctab}

\subsubsection*{Equipment}
8 envelopes or a tray \textbullet{} hand lens \textbullet{} ruler \textbullet{} pencil \textbullet{} sticky labels numbered 1 to 8

\subsubsection*{Method}
\begin{enumerate}[itemsep=2pt]
  \item Collect eight different leaves from the schoolyard. Label them 1 to 8 and write where each came from.
  \item Record the features of each leaf in Table 2.1.
  \item Choose a first question that splits the leaves into two groups. Keep asking questions until every leaf is on its own.
  \item Draw your key as a branching diagram. Give each leaf a name: its real name if you know it, or a made-up name if you do not.
  \item Swap leaves and keys with another group. Use their key to name each of their leaves, then check with them.
\end{enumerate}

% ─── RESULTS ─────────────────────────────────────────────────────────────────
\exppart{Results}{Amber}

\subsubsection*{Table 2.1: Features of eight schoolyard leaves \skr{PC6}}
\begin{ctab}
\begin{tabular}{|SW{1.4cm}|W{2.6cm}|C{2.2cm}|W{2.6cm}|C{2.4cm}|W{2.6cm}|}
  \hline\hd & \thead{Leaf} & \thead{Where found} & \thead{Length (cm)} & \thead{Margin (smooth / toothed)} & \thead{Leaflets (yes / no)} & \thead{Other feature} \\ \hline
  & \scaf{e.g.} & \scaf{Oval fence} & \scaf{11.5} & \scaf{Smooth} & \scaf{No} & \scaf{Curved, waxy} \\ \hline
  & 1 & & & & & \\ \hline
  & 2 & & & & & \\ \hline
  & 3 & & & & & \\ \hline
  & 4 & & & & & \\ \hline
  & 5 & & & & & \\ \hline
  & 6 & & & & & \\ \hline
  & 7 & & & & & \\ \hline
  & 8 & & & & & \\ \hline
\end{tabular}
\end{ctab}

\needspace{20\baselineskip}
\task Draw your dichotomous key as a branching diagram. \skr{CO2}
\drawbox{12cm}{}

\needspace{12\baselineskip}
\subsubsection*{Table 2.2: Testing another group's key \skr{EV3}}
\begin{ctab}
\begin{tabular}{|SW{1.4cm}|W{5.2cm}|W{5.2cm}|C{2.0cm}|}
  \hline\hd & \thead{Leaf} & \thead{Name from their key} & \thead{Correct name (from them)} & \thead{Correct?} \\ \hline
  & 1 & & & \\ \hline
  & 2 & & & \\ \hline
  & 3 & & & \\ \hline
  & 4 & & & \\ \hline
  & 5 & & & \\ \hline
  & 6 & & & \\ \hline
  & 7 & & & \\ \hline
  & 8 & & & \\ \hline
\end{tabular}
\end{ctab}

% ─── ANALYSIS ────────────────────────────────────────────────────────────────
\exppart{Analysis and evaluation}{Coral}

\task How many leaves did the other group identify correctly with your key? Which question caused a mistake? \skr{EV3}
\answerbox{4cm}

\needspace{10\baselineskip}
\task Rewrite that question so it could not be misread.
\answerbox{3cm}

\needspace{10\baselineskip}
\task A key is a \emph{model} for sorting living things. Give one thing your key does well, and one thing it could not do (for example, a leaf from a ninth plant). \skr{PA8}
\answerbox{4cm}

\needspace{10\baselineskip}
\task Why should a key use features you can see, rather than where the leaf was found or what season it is?
\answerbox{4cm}

\begin{yr8box}
\begin{keepext}
\exttask Field guides to Australian birds often use a key based on size, bill shape and colour. Explain why colour can be a poor feature for a key, using an example such as a young bird, a male and a female.
\answerbox{5cm}
\end{keepext}
\end{yr8box}

% ─── SKILL CHECK ─────────────────────────────────────────────────────────────
\exppart{Skill check}{NavyBlue}

\begin{skillcheck}
\scrow{PC5}{I collected specimens without harming living things.}
\scrow{PC6}{I recorded the features of each leaf carefully.}
\scrow{PA4}{I built a key that separates every specimen.}
\scrow{CO2}{I drew my key as a clear branching diagram.}
\scrow{EV3}{I found the question that caused errors and fixed it.}
\end{skillcheck}

\task Which skill will you focus on in Experiment 3? Explain what you will do differently.
\answerbox{5cm}
